% =============================================================================
% cap3_pipeline/02_fase1_ingestao.tex
% Seção 3.2 -- Fase I: Ingestão e pré-tratamento dos dados
%   Função da fase: homogeneizar a entrada heterogênea (vídeo H.264/H.265,
%   JPEG sequencial, mídia de trail cam) em um conjunto canônico de quadros
%   prontos para inferência, anexando metadados de proveniência por ponto.
%   Postura: descrever só o que é estritamente necessário para a operação
%   do pipeline; sem aprofundamentos típicos de processamento de imagens.
%   Texto corrido, impessoal, sem listas no corpo.
% =============================================================================
\section{Fase I --- Ingestão e triagem}
\label{sec:pipeline-fase1}

A Fase~I é responsável por converter a heterogeneidade de entrada descrita na Seção~\ref{sec:pipeline-visao-geral} em um \textbf{conjunto canônico de quadros} sobre o qual as fases seguintes possam operar sem precisar conhecer a origem específica de cada item. Sua função não é melhorar a qualidade das imagens em sentido fotográfico, mas garantir que o pipeline opere sobre uma representação uniforme e rastreável dos dados de campo, na qual cada quadro carrega consigo a informação necessária para que falhas possam ser diagnosticadas \emph{a posteriori}. Em projetos de \emph{camera trapping} em larga escala, esta etapa costuma absorver mais tempo de engenharia que as fases de inferência propriamente ditas, justamente porque dela depende a coerência de todo o restante~\cite{velez2022choosing,bruce2025largescale}.

\subsection*{Ingestão e extração de quadros}

A ingestão começa pela transferência das mídias gravadas em cartões microSD para um diretório de trabalho do servidor, organizado por ponto e por janela de coleta. No caso operacional fixado pela Seção~\ref{sec:decisao-arquitetural} --- IPCAM comercial doméstica adaptada em modo \emph{offline} ---, as mídias chegam como arquivos de vídeo em formatos \emph{H.264} ou \emph{H.265} (contêineres \texttt{.mp4} ou \texttt{.mkv}). Eventuais unidades experimentais paralelas (\emph{trail camera}, SBC com câmera, ESP32-CAM) entregam mídia em formatos distintos --- \emph{JPEG} estático, sequências numéricas, clipes proprietários --- tratada por adaptadores dedicados de ingestão que produzem a mesma representação lógica de quadro a jusante. Sobre os arquivos de vídeo, o pipeline executa uma \textbf{extração de quadros} (\emph{frame extraction}) com taxa de amostragem ajustada à dinâmica observada nos pontos: gatos comunitários permanecem junto aos potes de ração por janelas de segundos a alguns minutos, com movimentos relativamente lentos, o que permite amostrar em taxa significativamente inferior à do vídeo nativo --- da ordem de poucos quadros por segundo, configurável por ponto --- sem perda apreciável de informação relevante para detecção e re-ID. Essa decisão reduz drasticamente o volume de dados encaminhado às fases seguintes e elimina, já na entrada, quadros sucessivos quase idênticos.

Quando provenientes de unidades experimentais paralelas, imagens \emph{JPEG} estáticas são tratadas como quadros já amostrados e ingressam diretamente no estágio seguinte da Fase~I sem extração adicional. Ao fim deste sub-estágio, o pipeline trabalha sobre um \textbf{único formato lógico} --- um quadro de imagem com metadados associados --- independentemente da rota física pela qual o dado foi adquirido.

\subsection*{Filtragem de quadros não utilizáveis}

Antes de invocar qualquer modelo de visão, a Fase~I aplica duas filtragens \emph{baratas} (\emph{i.e.}, baseadas em estatísticas simples sobre histogramas) que eliminam quadros estruturalmente não utilizáveis e evitam que esses casos consumam custo computacional desnecessário nas fases seguintes. A primeira é o descarte de quadros \textbf{completamente escuros ou completamente saturados}, comuns em pontos onde a iluminação ativa por NIR (Seção~\ref{sec:acondicionamento-geometria}) gera quadros próximos ao módulo saturados e quadros distantes inteiramente pretos; essa filtragem é feita por inspeção simples da média e do desvio padrão de luminância do quadro. A segunda é o descarte de quadros sucessivos \textbf{quase idênticos} (deduplicação intra-evento), implementada por uma medida de diferença grosseira entre quadros adjacentes (por exemplo, soma de diferenças absolutas em uma versão reduzida do quadro). Essa deduplicação reconhece que um único evento de visitação ao pote, capturado em vídeo, pode gerar dezenas de quadros amostrados quase indistinguíveis, e que apenas alguns deles agregam informação efetiva para o pipeline.

Cabe registrar que ambas as filtragens da Fase~I são deliberadamente \emph{conservadoras}: o objetivo aqui é eliminar apenas os casos cuja inutilidade é estatisticamente evidente, repassando qualquer quadro com alguma chance de conter informação útil para o detector da Fase~II. Falsos positivos da Fase~I --- quadros não úteis que escapem desta filtragem inicial --- têm custo modesto, porque serão descartados na Fase~II; falsos negativos --- quadros úteis erroneamente descartados aqui --- são, em contraste, irrecuperáveis, e por isso a Fase~I privilegia recall sobre precisão.

\subsection*{Normalização e metadados por ponto}

Uma vez selecionados os quadros candidatos, a Fase~I executa uma \textbf{normalização leve} cujo objetivo é compensar diferenças sistemáticas entre pontos (\emph{e.g.}, diferentes regimes de exposição automática entre unidades, presença ou ausência de iluminação ativa por NIR, variações de \emph{white balance} sob luz mista) sem alterar o conteúdo semântico do quadro. Essa normalização restringe-se a operações simples e reversíveis --- equalização leve de histograma, ajuste de \emph{white balance} sob NIR e redimensionamento para a resolução de entrada esperada pelos modelos da Fase~II --- e segue diretrizes comuns em projetos de aquisição heterogênea com câmeras de armadilha~\cite{velez2022choosing,velasco2024smartcameratraps}. A intenção declarada é \emph{harmonizar} a entrada, e não \emph{embelezar} os quadros: artefatos de iluminação NIR, granulação noturna e variações angulares por ponto são informações úteis para a detecção e a classificação posteriores e devem ser preservados.

A cada quadro canônico que sai da Fase~I é anexado um conjunto mínimo de \textbf{metadados de proveniência}: identificador do ponto (P01--P10), tipologia da unidade que gerou a mídia (operacional ou experimental), \emph{timestamp} extraído do nome do arquivo ou dos cabeçalhos EXIF, e parâmetros de extração e normalização aplicados. Esses metadados não fazem parte da inferência em si, mas são instrumentais para o procedimento de avaliação descrito no Capítulo~\ref{cap:metodologia} e para o diagnóstico de falhas por ponto: sem eles, um decaimento de desempenho restrito a um determinado abrigo ou a uma determinada tipologia de hardware seria invisível à análise agregada. Ao final da Fase~I, portanto, o que segue para a Fase~II não é apenas uma coleção de imagens, mas uma coleção de quadros canônicos anotados com a informação mínima necessária para que o restante do pipeline opere e seja avaliável.
